% Nouveau EGA V (cote 25): archive cote 25, site page 72 (PDF page 73).
% typescript. Draft compared: none.
% AI-transcribed by Claude Opus 5.5 (claude-opus-5-5, adaptive thinking, effort high), Message Batches API, batch ega25b (msgbatch_01A7n5YuXubUPNK1zsepi8RG), 2026-11-13; not yet checked by eye.
\scanpage{25/72}
\mspage{67}

adt seulement l'adhérence de la sous-variété lisse formé des $\hand{\zeta}$ tels que
la dimension de $T\hand{\cap}L_{\hand{\zeta}}$ soit juste \uline{une} de plus que la dimension
``normal\unsure{è|e}'' n-m (n=dim T, m = codim L ). Sauf erreur, l'ensemble
(contenu dans l'ensemble supersingulier relatif) $V''$ introduit dans le par
16 (compléments) n'est autre que l'ensemble formé des couples
$(x,L_{\hand{\zeta}})$ tels que la dimension de $T_x\hand{\cap}L_{\hand{\zeta}}$ soit $\hand{\geqslant}$ n-m + 2\hand{)} , de sorte que
$\underline{X}$x $V^1 - V''$ \xout{\unsure{mx}}est \xout{xxxxxxxxx} lisse sur X, et sauf erreur est même \inserted{exactement} l'ensem
ble des points de $V^1$ lisses sur X.(La vér\xout{\unsure{i}}ification de ce point demanderai\unsure{t}
une étude de la filtration du \xout{xx} \unsure{shhéma|schéma} grassmanien suivant les \xout{xxxxx}
dimensions d'intersection des L variables avec une T fixe, sauf erreur on
trouve que le cran suivant de la filtration est formée exactement des
points non lisses du cran précédent - quand on définit la filtration
\margin{$|$} non seulement ensemblistement, mais aussi schématiquement, en utilisant
\margin{$|$} le lemme de la page 16 des compléments au par.16 \hand{\unsure{)}}. Cette étude forme-%
\margin{$|$} rait alors un des n$^{\text{o}}$s du paragraphe ``géométrique'' consacré aux grassmanien
\margin{$|$} nes \hand{)}. Si on définit de même \xout{V(k)} $V^{(k)}$ comme le sous-schéma de $X^{(m)}$
correspondant à dim $T_x\hand{\cap}L \geqslant$ n-m + k , on trouve par un calcul immédiat
que \xout{dimx$V^{(k)}$xxxxx\unsure{(k+1)(n+m)+k}x$^2$}
\[ \text{dim Grass}^m(P) - \text{dim } V^{(k)} = (k-1)(n-m) + k^2 \]
(du moins pour les restrictions raisonnables k $\hand{\leqslant}$ m, k $\hand{\leqslant}$r-m), sauf
erreur de calcul (NB ceci résulte plus généralement d\xout{xx}'un calcul des
dimensions des ``cellules'' qui interviennent dans la filtration des grass-%
mani\xout{x}ennes à laquelle j'ai fait allusion plus haut). Pour k = 2, on
trouve une différence de dimension $\geqslant 4$ , donc l'image de $V''$ dans Grass$^m$
est de codimension $\hand{\geqslant}$ 4, donc si on s'intéresse à ce qui se passe en dehors
d'ensembles \inserted[typed]{de Grass$^m$} de codimension $\geqslant$ 2, on peut \xout{xxx} oublier $V''$ .

D'autre part, dans $X_{\text{Grass}^m} - V''$ au dessus de Grass$^m$ , la situation
est celle du cas favorable envisa\hand{\unsure{gé}} dans les `` compléments au par 16'' ,

\mspage{72}
